\documentclass[11pt]{article}
\usepackage[margin=1in]{geometry}
\usepackage{amsmath,amssymb,amsthm,mathtools,microtype}
\usepackage[hidelinks]{hyperref}

\newtheorem{theorem}{Theorem}[section]
\newtheorem{lemma}[theorem]{Lemma}
\newtheorem{proposition}[theorem]{Proposition}
\newtheorem{corollary}[theorem]{Corollary}
\newtheorem{remark}[theorem]{Remark}
\newcommand{\A}{\mathbb A}
\renewcommand{\P}{\mathbb P}
\newcommand{\C}{\mathbb C}

\title{Hurwitz--LL Compactification of Rerooting}
\author{Repository working draft}
\date{\today}

\def\HurwitzLLPaper{1}

\begin{document}
\maketitle

\begin{abstract}
We compactify the rerooting relation for normalized marked-root pencils by
passing to marked admissible polynomial covers.  The resulting quotient-stack
description extends the generic degree-$N-2$ rerooting map, identifies its
boundary pullbacks, and compares the normalized Stein factors and conductor
square with the affine incidence constructions at arbitrary simultaneous root
collisions.  We also compute the degree $(N-2)N^{N-3}$ of the marked-zero-
fiber Lyashko--Looijenga morphism and the caustic and Maxwell divisor classes
on the root-boundary compactification.
\end{abstract}

\section{Setup, provenance, and review status}

\begin{quote}
\textbf{Audit notice (23 July 2026).}
The standalone \href{dvr-marking-audit.tex}{DVR marking audit} isolates the
missing global comparison.  The subsequent
\href{global-comparison-obstruction.tex}{global comparison obstruction}
shows that its first proposed arrow is false for the stated compactifications:
in degree five the target branch cross-ratio has leading term $x^3/y^2$, so
the admissible model requires a normalized graph blowup.  The abstract
quotient-stack map, finite-cover valuative lemma, and local root-collision
rings remain valid.  Claims identifying the full root-stable and
branch-stable compactifications should be read with the corrected graph
correspondence in that note.

The subsequent
\href{../../extended-geometry/BRANCH_SCALE_FAN.md}{branch-scale fan
experiment} computes the degree-six three-cluster chart.  Its radial part is
the predicted weighted braid fan, while a triple leading-value resonance
retains a second-scale cross-ratio and forces a further nested refinement.
The stable-target refinement in that chart is the normalized pullback of the
Kapranov four-point blowup of $\mathbb P^2$ giving $\overline M_{0,5}$; the
admissible source expansion and saturated-node normalization agree on the
generic equal-scale face and its first diagonal refinement.  All thirteen
radial scale types have verified degree-six source-component and saturated
node-monoid data.  All pairwise and triple Maxwell source-node rings and
their radial intersections are also verified; standard admissible-cover
deformation theory supplies their ambient toroidal gluing.  An exact central
Hurwitz enumeration finds two ambient classes, while the labelled
source-root cross-ratio selects the polynomial class as a reduced branch.
Thus the labelled coarse graphs agree in this chart.  Stack inertia, arising
from normalized label-preserving cover automorphisms, is trivial on every
radial type and is one diagonal $\mu_2$ at each pairwise or triple Maxwell
target node.  The matching algebraic candidate is the iterated second-root
stack along the four Maxwell boundary divisors.  Its face maps and
codimension-two inertia are explicit, and smooth tame-stack reconstruction
identifies it with the labelled ACV graph on this chart.  The separate
$(S_2)^3\rtimes S_3$ quotient is explicit.  Gluing under the full root-label
action is canonical on the global labelled normalized graph closure.  The
stable-target graph on every collision type is the normalized wonderful
pullback of the branch-diagonal building set.  The companion
\href{../../extended-geometry/SOURCE_VERTEX_RIGIDITY.md}{source-vertex
rigidity theorem} proves that two target-flag fibers and one third-fiber
equation uniquely reconstruct each rational source-component map.  The
general radial source theorem constructs those components and their
saturated node partitions on every ordered first-scale stratum for arbitrary
cluster multiplicities.  Polynomial monodromy subforests determine the
source dual trees, component degrees, and node indices on every nested
simple-branch resonance stratum.  Since the fully marked admissible-cover
branch morphism is finite, finite normalization proves that the wonderful
graph is already the complete coarse polynomial source graph.  Corrected H2
and coarse H3 are therefore unconditional.  Monodromy centralizers also
give the full-chain radial formula: equal cluster multiplicities have
trivial inertia, whereas unequal multiplicities can retain a finite
subgroup.  They also compute generic simple-resonance inertia from anchors
and node partitions.  The companion
\href{../../extended-geometry/RECURSIVE_RESONANCE_ATLAS.md}{recursive
resonance atlas theorem} gives normalized-initial-form formulas for source
flag divisors in framed residue coordinates, contraction descent by affine
composition, and simultaneous centralizer characters.  Its selected-factor
log-\'etale comparison supplies global coverage, overlap descent, and
identification with the normalized admissible-cover stack.  Thus
H1-COARSE, H1-STACK, H2, and coarse H3 are complete.
\end{quote}

Let $H$ be a normalized degree-$N$ seed with an exact double root at zero and
otherwise simple zero-fiber points on the dense open under consideration.
The construction-independent
\href{../../verified/TANGENT_MAP_CORE.md}{tangent-map core} supplies the
common normal form
\[
 (W,s)\longmapsto\bigl(s,Ws-H(W)\bigr).
\]
For
\[
 P_{H,s}(W)=H(W)-sW,
\]
its critical restriction, a point $r$ has $s=H'(r)$ and critical value
$H(r)-rH'(r)$.  Thus the discriminant normalization is
\[
 r\longmapsto \bigl(H'(r),rH'(r)-H(r)\bigr).
\]
The $N-2$ nonzero simple roots give the generic rerootings; collision strata
are treated in the marked admissible-cover closure.

The generic affine marking is supplied by the unique unramified root sheet
meeting the regular-reconstruction open.  On the marked admissible-cover
graph the selected root extends tautologically: if
$\Gamma_N^{\rm lab}$ is the fully labelled normalized graph, the
marked-to-unmarked map
\[
 [\Gamma_N^{\rm lab}/S_{N-3}]
 \longrightarrow
 [\Gamma_N^{\rm lab}/S_{N-2}]
\]
is finite, representable, and etale of degree $N-2$.  The normalized Stein
contractions, completed root-cover charts, and conductor squares are also
identified here at arbitrary simultaneous multicluster collisions.  After
coarse contraction, the finite birational closure of the distinguished
regular-reconstruction component is an isomorphism over the normal marked
coarse compactification, so the mark has a unique point over every DVR limit.

The admissible-cover technology cited below from Abramovich--Corti--Vistoli
and Deopurkar supplies an appropriate proper compactification framework.  It
does not externally review the repository-specific restricted LL-degree,
boundary-pullback, Picard, caustic, or Maxwell calculations.  No external
review of those calculations is recorded, and the statements below should be
read with that provenance distinction explicit.

\section{The scheme-theoretic decorated normalization}

For a weighted primitive $H$, let $D_H$ be the reduced discriminant of
$H(W)-sW+t$ and let
\[
 \nu_H:\A^1_r\longrightarrow D_H,
 \qquad r\longmapsto(H'(r),rH'(r)-H(r))                 \tag{6.1}
\]
be its affine normalization.  Write
$R_H=k[H'(r),rH'(r)-H(r)]\subset k[r]$.

\begin{proposition}[full Fitting divisor]
\label{prop:full-fitting}
There is an equality of ideal sheaves on the normalization
\[
 \operatorname{Fitt}_0\Omega_{\widetilde D_H/D_H}=(H''(r)). \tag{6.2}
\]
This ideal, including every root multiplicity, is preserved by polynomial
left--right equivalence and by adjoining identity variables.
\end{proposition}

\begin{proof}
The images in $k[r]\,dr$ of the two generators of $R_H$ are
\[
 dH'(r)=H''(r)dr,
 \qquad d(rH'(r)-H(r))=rH''(r)dr.
\]
Thus
$\Omega_{k[r]/R_H}\simeq k[r]/(H'')\,dr$, proving (6.2).
The stable boundary theorem identifies the finite normalization morphisms on
the intrinsic discriminant vertex away from the second vertex, so it
identifies their relative differential modules and Fitting ideals.  Under
stabilization one has
\[
 \Omega_{(\widetilde D_H\times\A^q)/(D_H\times\A^q)}
 \simeq\operatorname{pr}_1^*\Omega_{\widetilde D_H/D_H}.
\]
Zeroth Fitting ideals commute with this flat base change.  Equality of the
ideal sheaves preserves their valuations at every height-one point, hence
preserves the multiplicities in (6.2), not only the reduced support.
\end{proof}

If the support of (6.2) contains two points, the units of
$k[r,\xi,\xi^{-1},z_1,\ldots,z_q]$ and subtraction of the two corresponding
principal-prime equations force the induced coordinate change to have the
form $r\mapsto ar+b$, $a\ne0$.  Proposition~\ref{prop:full-fitting} then says
that the full effective Hessian divisor is carried affinely, with its
coefficients.

The self-intersection of (6.1) is also intrinsic.  In two normalization
coordinates put
\[
 D_s(r,u)=\frac{H'(r)-H'(u)}{r-u},
\]
\[
 D_t(r,u)=
 \frac{rH'(r)-H(r)-uH'(u)+H(u)}{r-u}.
\]
The genuine off-diagonal closure is the scheme
\[
 \Gamma_H^{\rm off}
 =V\bigl((D_s,D_t):(r-u)^\infty\bigr).                 \tag{6.3}
\]
The saturation in (6.3) is necessary: on the diagonal the divided equations
specialize to $H''(r)$ and $rH''(r)$ and otherwise retain isolated cusp
points.  The involution $(r,u)\leftrightarrow(u,r)$ is free on the ordinary
off-diagonal locus; its characteristic-zero finite quotient exists in
general and retains fixed diagonal limits in degenerations.  On the ordinary
locus that quotient records the two normalization branches paired at each node.  An
isomorphism of normalization morphisms preserves the fiber product, its
diagonal, and the saturated complement, while stabilization takes their
product with affine space.  Hence this node-pairing scheme is stable.
In the exact extractor, ordinary-node status is certified by the Jacobian
determinant
\[
 \det\frac{\partial(D_s,D_t)}{\partial(r,u)}
\]
being a unit on (6.3), rather than inferred only from the number of projected
branch points.  Passing to
$\sigma=r+u,\pi=ru$ gives the scheme quotient by the free involution; the
implementation verifies that its equations pull back into the saturated
ordered ideal.

Finally, let
\[
 \mathfrak c_H=\operatorname{Ann}_{\mathcal O_{D_H}}
 (\nu_{H*}\mathcal O_{\widetilde D_H}/\mathcal O_{D_H}) \tag{6.4}
\]
be the conductor.  It is functorial for the same reason.  On the open locus
of ordinary cusps and nodes, if $P_H(r)$ cuts out all normalization branches
appearing in (6.3), then the local models
$k[[t^2,t^3]]\subset k[[t]]$ and
$k[[x,y]]/(xy)\subset k[[x]]\oplus k[[y]]$ give
\[
 \mathfrak c_Hk[r]=(H''(r)^2P_H(r)).                   \tag{6.5}
\]
For arbitrary plane-curve singularities, let $f_H(S,T)=0$ be the primitive
implicit equation.  Plane-curve adjunction gives the uniform generator
\[
 \mathfrak c_Hk[r]=
 \left(
 \frac{(\partial f_H/\partial T)
 (H'(r),rH'(r)-H(r))}{H''(r)}
 \right).                                             \tag{6.5a}
\]
This quotient is polynomial because (6.1) is the normalization.  The exact
extractor uses (6.5a) without an ordinary-singularity hypothesis and checks
(6.5) independently whenever the cusp/node criterion holds.
The invariant retains the finite conductor morphism
\[
 \operatorname{Spec}(k[r]/\mathfrak c_Hk[r])
 \longrightarrow
 \operatorname{Spec}(R_H/(R_H\cap\mathfrak c_Hk[r])), \tag{6.6}
\]
not only its upstairs support.  It records the two branches over a node and
the single thick branch over a cusp.  Together, (6.2)--(6.6), the point at infinity, and the labeled inverse image
of the intersection with the second target boundary form the decorated
normalization.  The full second-boundary center divisor upstairs is
$\gcd(H,H')$; a point label is inferred only when its support is uniquely
certified.  Every part is therefore preserved under stable polynomial
left--right equivalence.

For the split quartics, the boundary label on the normalization is also
checked intrinsically.  If zero is the unique multiple primitive root and
$H(r)=h_2r^2+O(r^3)$, the normalized discriminant boundary chart
\[
 B=\frac{H'(r)}C,\qquad
 A=\frac{rH'(r)-H(r)}{cC^2},\qquad r=Ck
\]
specializes to
\[
 B=2h_2k,\qquad A=\frac{h_2}{c}k^2.
\]
A finite specialization centered at $r=\rho$ forces
$H(\rho)=H'(\rho)=0$, hence $\rho=0$.  Thus the second-boundary intersection
canonically marks $r=0$; it is not an auxiliary normalization convention.

For completeness, this decoration gives an actual generic-dimension theorem,
not only a parameter count.  In the normalized degree-$N$ seed space, an
isomorphism preserving zero and infinity is $r\mapsto ar$.  Equality of
Fitting divisors then has the form $G''(r)=cH''(ar)$.  Since the normalized
seeds vanish together with their first derivatives at zero, integration gives
\[
 G(r)=\frac{c}{a^2}H(ar).
\]
The condition $G(1)=0$ forces $H(a)=0$, leaving at most $N-2$ generic
choices, and $G'(1)=-1$ determines $c$.  Thus the decorated-normalization
moduli map is generically finite and its image has dimension $N-3$.


\end{document}
